% Generated from locale/tr/content/intuitionistic-logic/introduction/natural-deduction.tex; do not hand-edit.


\olfileid[tr]{int}{int}{ntd}

\olsection{Doğal Tümdengelim}

$\FalseCl$ kuralları bulunmayan doğal tümdengelim, sezgici mantık için
standart bir !!{derivation} sistemidir. Kuralları burada yeniden veriyor ve
güdülerini BHK yorumunu kullanarak açıklıyoruz. Her durumda kuralı, öncüllerin
inşaları varsa sonucun da bir inşası bulunduğu sonucuna varmamıza izin veren
bir kural olarak düşünebiliriz.

Doğal tümdengelim !!{derivation}s{}inde kapatılmamış varsayımlar bulunduğundan,
böyle !!a{derivation}{}i---sözgelimi $!A$'yı kapatılmamış~$\Gamma$
!!{undischarged} varsayımlarından türeten birini---tüm $!B \in \Gamma$'ların
inşalarını~$!A$'nın bir inşasına dönüştüren bir fonksiyon olarak görmeliyiz.
!!a{derivation},~$!A$'yı hiçbir !!{undischarged} varsayıma dayanmadan veriyorsa
BHK yorumu anlamında~$!A$'nın bir inşası vardır. Bununla
birlikte tartışma sırasında gerekmediğinde~$\Gamma$'yı göstermeyeceğiz.

Tek başına bir~$!A$ varsayımı,~$!A$'nın !!a{derivation}{}idir; buradaki
!!{undischarged} varsayım~$!A$'dır. Bu, BHK yorumuyla uyuşur: inşalar üzerindeki
özdeşlik fonksiyonu,~$!A$'nın her inşasını~$!A$'nın bir inşasına dönüştürür.

\subsection{Ve Bağlacı}

\begin{defish}
\AxiomC{$!A$}
\AxiomC{$!B$}
\RightLabel{\Intro{\land}}
\BinaryInfC{$!A \land !B$}
\DisplayProof
\hfill
\begin{tabular}{l}
\AxiomC{$!A \land !B$}
\RightLabel{\Elim{\land}}
\UnaryInfC{$!A$}
\DisplayProof
\\[3ex]
\AxiomC{$!A \land !B$}
\RightLabel{\Elim{\land}}
\UnaryInfC{$!B$}
\DisplayProof
\end{tabular}
\end{defish}

Birer $N_1$, $N_2$ inşasına sahip olduğumuzu varsayın; bunlar sırasıyla
$!A_1$ ile $!A_2$'nin inşalarıdır. O zaman $!A_1 \land !A_2$'nin de bir
inşasına, yani $\tuple{N_1,
N_2}$ çiftine sahibiz.

BHK yorumuna göre $!A_1 \land !A_2$'nin bir inşası bir $\tuple{N_1, N_2}$
çiftidir. Öyleyse böyle bir çiftimiz olduğunu varsayın. Bu durumda her bir
bileşenin de bir inşasına sahibiz: $N_1$,~$!A_1$'in; $N_2$ ise $!A_2$'nin
bir inşasıdır.

\subsection{Koşul}

\begin{defish}
  \AxiomC{$\Discharge{!A}{u}$}
  \DeduceC{$!B$}
  \DischargeRule{$\Intro{\lif}$}{u}
  \UnaryInfC{$!A \lif !B$}
  \DisplayProof
\hfill
  \AxiomC{$!A \lif !B$}
  \AxiomC{$!A$}
  \RightLabel{$\Elim{\lif}$}
  \BinaryInfC{$!B$}
  \DisplayProof
\end{defish}

$!B$'nin kapatılmamış~$!A$ varsayımından !!a{derivation}{}ine ve bu varsayımın
!!{undischarged} kaldığı bir türetime sahipsek bir~$f$ fonksiyonu vardır; bu
fonksiyon~$!A$'nın inşalarını~$!B$'nin inşalarına dönüştürür. Aynı fonksiyon
$!A \lif !B$'nin bir inşasıdır. Dolayısıyla
$\Intro\lif$ kuralının öncülünün~$!A$'nın bir inşasına bağlı bir inşası varsa,
sonuç $!A \lif !B$'nin de bir inşası vardır.

Öte yandan, bir~$N$ inşası bulunduğunu, bunun~$!A$'nın inşası olduğunu ve
ayrıca bir~$f$ inşası bulunduğunu, bunun da $!A \lif !B$'nin inşası olduğunu
varsayın. $!A \lif !B$'nin inşası,
$!A$'nın inşalarını~$!B$'nin
inşalarına dönüştüren bir fonksiyondur. Öyleyse $f(N)$,~$!B$'nin bir inşasıdır;
yani $\Elim\lif$ kuralının sonucunun bir inşası vardır.

\subsection{Veya Bağlacı}

\begin{defish}
\begin{tabular}{l}
\AxiomC{$!A$}
\RightLabel{\Intro{\lor}}
\UnaryInfC{$!A \lor !B$}
\DisplayProof
\\[3ex]
\AxiomC{$!B$}
\RightLabel{\Intro{\lor}}
\UnaryInfC{$!A \lor !B$}
\DisplayProof
\end{tabular}
\hfill
\AxiomC{$!A \lor !B$}
\AxiomC{$\Discharge{!A}{n}$}
\DeduceC{$!C$}
\AxiomC{$\Discharge{!B}{n}$}
\DeduceC{$!C$}
\DischargeRule{\Elim{\lor}}{n}
\TrinaryInfC{$!C$}
\DisplayProof
\end{defish}

Bir~$N_i$ inşasına sahipsek ve bu~$!A_i$'nin inşasıysa, onu bir
$\tuple{i, N_i}$ inşasına, yani~$!A_1 \lor !A_2$'nin inşasına
dönüştürebiliriz. Öte yandan, $!A_1 \lor !A_2$'nin bir inşasına, yani bir
$\tuple{i, N_i}$ çiftine sahip olduğumuzu varsayın; burada $N_i$,~$!A_i$'nin
inşasıdır. Ayrıca $f_1$, $f_2$ fonksiyonlarının sırasıyla $!A_1$, $!A_2$
inşalarını~$!C$ inşalarına dönüştürdüğünü varsayın. O zaman $f_i(N_i)$,
~$!C$'nin, yani~$\Elim\lor$ kuralının sonucunun bir inşasıdır.

\subsection{Saçmalık}

\begin{defish}
  \AxiomC{$\lfalse$}
  \RightLabel{\FalseInt}
  \UnaryInfC{$!A$}
  \DisplayProof
\end{defish}

$\lfalse$'ın kapatılmamış~$!B_1$, \dots, $!B_n$ varsayımlarından
!!a{derivation}{}ine ve bu varsayımların !!{undischarged} kaldığı bir türetime
sahipsek, bir $f(M_1, \dots, M_n)$ fonksiyonu vardır; bu fonksiyon $!B_1$,
\dots, $!B_n$'nin inşalarını $\lfalse$'ın bir inşasına dönüştürür.
$\lfalse$'ın hiçbir inşası olmadığından, $!B_1$, \dots, $!B_n$'nin tümünün
inşaları da bulunamaz. Dolayısıyla $f$ şu özelliğe de sahiptir:
\emph{eğer} $M_1$, \dots, $M_n$ sırasıyla $!B_1$, \dots, $!B_n$'nin
inşalarıysa, \emph{o zaman} $f(M_1, \dots, M_n)$,~$!A$'nın bir inşasıdır.

\subsection{$\lnot$ için Kurallar}

$\lnot !A$, $!A \lif \lfalse$ olarak tanımlandığından, kesin konuşmak
gerekirse~$\lnot$ için kurallara ihtiyacımız yoktur. Yine de kurallarımız
olsaydı şöyle görünürlerdi:

\begin{defish}
\AxiomC{$\Discharge{!A}{n}$}
\noLine
\DeduceC{$\lfalse$}
\DischargeRule{\Intro{\lnot}}{n}
\UnaryInfC{$\lnot !A$}
\DisplayProof
\hfill
\AxiomC{$\lnot !A$}
\AxiomC{$!A$}
\RightLabel{\Elim{\lnot}}
\BinaryInfC{$\lfalse$}
\DisplayProof
\end{defish}


\subsection{\usetoken{P}{derivation} Örnekleri}

\begin{enumerate}
\item $\Proves !A \lif (\lnot !A \lif \lfalse)$, yani $\Proves !A
  \lif ((!A \lif \lfalse) \lif \lfalse)$
  \begin{prooftree}
    \AxiomC{$\Discharge{!A}{2}$}
    \AxiomC{$\Discharge{!A \lif \lfalse}{1}$}
    \RightLabel{\Elim\lif}
    \BinaryInfC{$\lfalse$}
    \DischargeRule{\Intro\lif}{1}
    \UnaryInfC{$(!A \lif \lfalse)\lif \lfalse$}
    \DischargeRule{\Intro\lif}{2}
    \UnaryInfC{$!A \lif (!A \lif \lfalse) \lif \lfalse$}
  \end{prooftree}

\item $\Proves ((!A \land !B) \lif !C) \lif (!A \lif (!B \lif !C))$
  \begin{prooftree}
    \AxiomC{$\Discharge{(!A \land !B) \lif !C}{3}$}
    \AxiomC{$\Discharge{!A}{2}$}
    \AxiomC{$\Discharge{!B}{1}$}
    \RightLabel{\Intro\land}
    \BinaryInfC{$!A \land !B$}
    \RightLabel{\Elim\lif}
    \BinaryInfC{$!C$}
    \DischargeRule{\Intro\lif}{1}
    \UnaryInfC{$!B \lif !C$}
    \DischargeRule{\Intro\lif}{2}
    \UnaryInfC{$!A \lif (!B \lif !C)$}
    \DischargeRule{\Intro\lif}{3}
    \UnaryInfC{$((!A \land !B) \lif !C) \lif (!A \lif (!B \lif !C))$}
  \end{prooftree}

\item $\Proves \lnot(!A \land \lnot !A)$, yani $\Proves (!A \land (!A
  \lif \lfalse))\lif \lfalse$
  \begin{prooftree}
    \AxiomC{$\Discharge{!A \land (!A \lif \lfalse)}{1}$}
    \RightLabel{\Elim\land}
    \UnaryInfC{$!A \lif \lfalse$}
    \AxiomC{$\Discharge{!A \land (!A \lif \lfalse)}{1}$}
    \RightLabel{\Elim\land}
    \UnaryInfC{$!A$}
    \RightLabel{\Elim\lif}
    \BinaryInfC{$\lfalse$}
    \DischargeRule{\Intro\lif}{1}
    \UnaryInfC{$(!A \land (!A \lif \lfalse))\lif \lfalse$}
  \end{prooftree}

\item $\Proves \lnot\lnot(!A \lor \lnot !A)$, yani $\Proves ((!A \lor
  (!A \lif \lfalse))\lif \lfalse)\lif \lfalse$
  \begin{prooftree}\hskip-3em
    \AxiomC{$\Discharge{(!A \lor (!A \lif \lfalse))\lif \lfalse}{2}$}
    \AxiomC{$\Discharge{(!A \lor (!A \lif \lfalse))\lif \lfalse}{2}$}
    \AxiomC{$\Discharge{!A}{1}$}
    \RightLabel{\Intro\lor}
    \UnaryInfC{$!A \lor (!A \lif \lfalse)$}
    \RightLabel{\Elim\lif}
    \BinaryInfC{$\lfalse$}
    \DischargeRule{\Intro\lif}{1}
    \UnaryInfC{$!A \lif \lfalse$}
    \RightLabel{\Intro\lor}
    \UnaryInfC{$!A \lor (!A \lif \lfalse)$}
    \RightLabel{\Elim\lif}
    \insertBetweenHyps{\hskip -4em}
    \BinaryInfC{$\lfalse$}
    \DischargeRule{\Intro\lif}{2}
    \UnaryInfC{$((!A \lor (!A \lif \lfalse))\lif \lfalse)\lif \lfalse$}
  \end{prooftree}
\end{enumerate}

\begin{prop}
  Sezgici mantık için doğal tümdengelimde $\Gamma \Proves !A$ ise klasik
  mantıkta da $\Gamma \Proves !A$ olur. Özellikle $!A$ sezgici bir teoremse
  klasik bir teoremdir.
\end{prop}

\begin{proof}
  Her doğal tümdengelim kuralı aynı zamanda klasik doğal tümdengelimin de bir
  kuralıdır; dolayısıyla sezgici mantıktaki her !!{derivation}, klasik mantıkta
  da !!a{derivation}{}dir.
\end{proof}

\begin{prob}
  Sezgici mantıkta !!{derivation} vererek aşağıdaki !!{formula}s{}i türetin:
  \begin{enumerate}
    \item $(\lnot !A \lor !B) \lif (!A \lif !B)$
    \item $\lnot\lnot\lnot !A \lif \lnot !A$
    \item $\lnot\lnot (!A \land !B) \liff (\lnot\lnot !A\land \lnot \lnot !B)$
    \item $\lnot (!A \lor !B) \liff (\lnot !A \land \lnot !B)$
    \item $(\lnot !A \lor \lnot !B) \lif \lnot (!A \land !B)$
    \item $\lnot\lnot ( !A \land !B) \lif  (\lnot\lnot !A \lor
    \lnot\lnot !B)$
  \end{enumerate}
\end{prob}

